\section{Resultados del análisis técnico}
\subsection{Opciones de despliegue y responsabilidades}
El cuadro~\ref{tab:plataformas} compara decisiones operativas. Binder 2.0 describe la publicación de entornos interactivos a partir de repositorios \cite{binderpaper}; el servicio público mybinder.org establece límites para sesiones efímeras \cite{binder}. JupyterHub permite administrar entornos por usuario y autenticación bajo control del operador \cite{jupyterhub}. Colab ofrece cuadernos gestionados con disponibilidad y límites variables \cite{colab}. Codespaces añade un entorno de desarrollo asociado al repositorio, con consumo de cómputo y almacenamiento facturable \cite{billing}.

\begin{table}[htbp]\centering\small
\caption{Comparación operativa según documentación de las plataformas consultada en octubre de 2026.}\label{tab:plataformas}
\begin{tabularx}{160mm}{@{}p{24mm}XXX@{}}\toprule
Opción & Administración & Conservación del trabajo & Condición a planificar \\
Codespaces & Servicio gestionado; configuración del repositorio & Espacio persistente hasta su eliminación; cambios a Git & Cuentas, cuotas, cómputo y almacenamiento \\
Binder público & Servicio comunitario; entorno desde repositorio & Sesión efímera; exportar cambios & Límites de sesión y capacidad compartida \\
JupyterHub & Operador institucional o proveedor & Depende del almacenamiento configurado & Autenticación, capacidad, respaldo y soporte \\
Colab & Cuadernos gestionados & Cuaderno guardado; estado de ejecución temporal & Límites variables y dependencias adicionales \\
\bottomrule\end{tabularx}
\end{table}


La comparación permite una decisión concreta: un ejercicio breve que solo necesita ejecutar un cuaderno tiene requisitos diferentes de un curso que enseña ramas, revisión de cambios y desarrollo de software. Codespaces encaja en este segundo objetivo si hay cuentas, conectividad y presupuesto; un JupyterHub institucional puede concentrar la administración cuando existe soporte local. Binder facilita demostraciones compartibles, mientras Colab ofrece una entrada rápida a cuadernos. La elección debe incluir el procedimiento para exportar el trabajo del alumno y conservarlo al terminar.

\subsection{Contexto docente e inventario de evidencia}
Los ejercicios públicos asociados a SILAFAE permiten examinar materiales de formación sin deducir de ellos tiempos de incorporación históricos. La 47.ª CERN School of Computing se celebró en Santiago del 11 al 24 de enero de 2026 \cite{csc}; su página de docentes incluye a Arturo Sánchez \cite{csclecturers}. Estos registros sitúan el contexto formativo. El cuadro~\ref{tab:evidencia} explicita qué evidencia sustenta cada afirmación del ensayo.

\begin{table}[tbp]
\caption{Evidencia de la trayectoria docente y alcance de cada registro.}\label{tab:evidencia}
\begin{tabularx}{\linewidth}{@{}>{\RaggedRight\arraybackslash}p{34mm}YY@{}}
\toprule
\cabecera{Registro} & \cabecera{Evidencia} & \cabecera{Alcance} \\
\midrule
\fila{SILAFAE, noviembre 2024} & Dos repositorios con SHA fechado & Contenido y recetas; no commit por alumno \\
\addlinespace
\fila{CSC, enero 2026} & Programa con responsables y horarios & Secuencia prevista; no tiempo real de preparación \\
\addlinespace
\fila{Selección CSC} & 70 plazas en diciembre; 69 seleccionados en enero & Cortes distintos; no asistencia a una sesión \\
\addlinespace
\fila{Ejercicios CSC} & SHA previo al curso y corrección del 12 de enero & Cambios en un ejercicio; no incidente operativo \\
\addlinespace
\fila{Evaluación por sesión} & Protocolo y plantillas de registro & Requiere consumo, incidencias y observaciones \\
\bottomrule
\end{tabularx}
\end{table}


El resultado técnico de esta lectura es separar tres objetos que suelen confundirse: un repositorio disponible, un entorno que ejecuta una tarea y una experiencia docente evaluada. Cada uno necesita una comprobación diferente. El repositorio permite inspección; una prueba limpia verifica ejecución; un registro por sesión permite estimar tiempos, incidencias y recursos. La arquitectura organiza esas comprobaciones sin sustituir una por otra.
